\documentclass[12pt,a4paper]{article}
\usepackage[utf8]{inputenc}
\usepackage[T1]{fontenc}
\usepackage{amsmath,amssymb,amsthm,mathrsfs,mathtools}
\usepackage[margin=1in]{geometry}
\usepackage{hyperref}
\hypersetup{colorlinks=true,linkcolor=blue,citecolor=blue,urlcolor=blue}
\usepackage{booktabs}

\theoremstyle{plain}
\newtheorem{theorem}{Theorem}[section]
\newtheorem{proposition}[theorem]{Proposition}
\newtheorem{lemma}[theorem]{Lemma}
\newtheorem{corollary}[theorem]{Corollary}
\theoremstyle{definition}
\newtheorem{definition}[theorem]{Definition}
\theoremstyle{remark}
\newtheorem{remark}[theorem]{Remark}

\newcommand{\R}{\mathbb{R}}
\newcommand{\C}{\mathbb{C}}
\newcommand{\dd}{\mathrm{d}}

\title{Principal-Series Blocks on the Celestial Sphere:\\[4pt]
Conical Reduction, Shadow Symmetry, and Exact Moments}

\author{Daniel Toupin\\[4pt]
\textit{Golden Physics Project}\\[2pt]
\texttt{goldenphysics.org}}

\date{August 2026 --- v2, superseding \texttt{kinematic\_block\_v1.1}}

\begin{document}
\maketitle

\begin{abstract}
\noindent
This paper collects the special-function structure of scalar
principal-series blocks on the celestial sphere, and re-anchors it as a
toolkit for the partial-wave decomposition of unitarity cuts. Six
results are stated and proved, each verified to at least 25 digits.
The chiral block $k_h(z)=z^h\,{}_2F_1(h,h;2h;z)$ reduces exactly to a
Legendre function of the second kind, $k_h(z)=c(h)\,Q_{h-1}(2/z-1)$,
in the off-cut branch appropriate to argument greater than one; the
branch choice is not cosmetic, and the identity fails by roughly $50\%$
in the on-cut convention. The shadow involution $\Delta\mapsto
2-\Delta$ becomes the classical Legendre degree symmetry
$\nu\mapsto-\nu-1$, whose connection formula
$Q_{-\nu-1}-Q_\nu=i\pi\tanh(\pi\tau)P_\nu$ exhibits the shadow-odd part
of the block as a Mehler--Fock kernel. The reduction prefactor has the
exact modulus $|c(\lambda)|^2=(2\lambda/\pi)\coth(\pi\lambda/2)$. On
physical kinematics the full block is bounded in $\lambda$, with
envelope $|1-z|^{-1/2}$, since the chiral exponential rates cancel by
conjugate oddness. We give a short closed-form proof, previously only a
numerical observation, of the first moment
\[
\frac{1}{2\pi}\int_{\R} P(\lambda)\,
\mathrm{Re}\,\psi\!\left(\tfrac12+\tfrac{i\lambda}{2}\right)\dd\lambda
\;=\;\tfrac18+\tfrac14\,\psi(\tfrac12),
\qquad
P(\lambda)=\frac{\pi\lambda}{\sinh(\pi\lambda)} ,
\]
by Fourier inversion followed by an elementary substitution.
Two corrections to the earlier draft are recorded. The value $1/8$ in
that moment and the one-loop ladder moment $\mathcal M_1=1/8$ use
different integration conventions (full line against half line), so
their agreement is a proved coincidence of two separate evaluations,
with no derivation linking them. And the vanishing of the Mellin bridge
$(1-2^{-s})\Gamma(s)\zeta(s)$ at nontrivial zeros carries no
information, since $\zeta(s)$ appears in it as a factor.
\end{abstract}

\medskip
\noindent\textbf{Keywords:} celestial holography, conformal blocks,
principal series, Legendre functions, Mehler--Fock transform.

\tableofcontents
\newpage

% =====================================================================
\section{Purpose and scope}
% =====================================================================

A companion paper \cite{ToupinLoops} derives the two-particle unitarity
cut in celestial variables and shows that it is uniform on the round
celestial sphere,
\begin{equation}
\dd\Pi_2=\frac{\dd^2z}{8\pi^2(1+|z|^2)^2},
\qquad
z_6=-\frac{1}{\bar z_5},
\label{eq:lips}
\end{equation}
with the two cut legs at antipodal points. To decompose that cut into
$\mathrm{SL}(2,\C)$ partial waves, one needs the scalar
principal-series block as an explicit function, its behaviour under the
shadow map, its size at large spectral parameter, and its moments
against the weight $P(\lambda)$ that the same companion derives from
phase space. This paper supplies those four things.

The earlier draft framed this material as the ``kinematic factor'' of a
factorization in which $P(\lambda)$ served as a replacement for the
Feynman loop measure. That framing is withdrawn in
\cite{ToupinSpectral}. The mathematics below is untouched by the
withdrawal, since none of it depends on the interpretation. What
changes is the use to which it is put.

\paragraph{Conventions.}
$h=\tfrac12(1+i\lambda)$ with $\lambda\in\R$ on the principal series;
$\nu=h-1=-\tfrac12+i\tau$, $\tau=\lambda/2$;
$x(z)=2/z-1$; $\psi$ is the digamma function. All numerical claims come
from the archived script \texttt{verify\_blocks\_v2.py}.

% =====================================================================
\section{Conical reduction, with its branch}
% =====================================================================

\begin{theorem}[Conical reduction]
\label{thm:conical}
For $h\notin\{0,-\tfrac12,-1,\dots\}$ and $z\in(0,1)$,
\begin{equation}
k_h(z)\;:=\;z^h\,{}_2F_1(h,h;2h;z)
\;=\;c(h)\,Q_{h-1}\!\left(\frac{2}{z}-1\right),
\qquad
c(h)=\frac{2\,\Gamma(2h)}{\Gamma(h)^2},
\label{eq:conical}
\end{equation}
where $Q_\nu$ is taken in the branch analytic on $(1,\infty)$ (the
``type 3'' or off-cut Legendre function), which is the branch selected
by the argument $x=2/z-1>1$.
\end{theorem}

\begin{proof}
Quadratic transformation of ${}_2F_1$ into Legendre normal form,
followed by identification of the recessive solution at $x\to\infty$.
Verified at four points, including complex $h$ on and off the principal
series, to relative error below $3\times10^{-41}$.
\end{proof}

\begin{remark}[The branch is load-bearing]
\label{rem:branch}
Evaluated in the on-cut convention (Legendre $Q$ with cut along
$(-\infty,1]$), \eqref{eq:conical} fails by $47\%$ at
$h=\tfrac12+0.7i,\,z=0.3$, and by $98\%$ at $h=0.8,\,z=0.25$. The
earlier draft did not record the convention. Any reader reproducing the
identity in a computer algebra system whose default is the on-cut
branch will see it fail, so we state it here.
\end{remark}

% =====================================================================
\section{Shadow symmetry as Legendre degree symmetry}
% =====================================================================

\begin{theorem}[Shadow involution and the Mehler--Fock kernel]
\label{thm:shadow}
Under $\Delta\mapsto2-\Delta$, equivalently $h\mapsto1-h$ and
$\nu\mapsto-\nu-1$, the reduction \eqref{eq:conical} maps $Q_\nu$ to
$Q_{-\nu-1}$, and for $\nu=-\tfrac12+i\tau$, $x>1$,
\begin{equation}
Q_{-\nu-1}(x)-Q_\nu(x)\;=\;i\pi\tanh(\pi\tau)\,P_\nu(x).
\label{eq:connection}
\end{equation}
The shadow-odd part of the block is therefore the conical function
$P_{-1/2+i\tau}$, the kernel of the Mehler--Fock transform.
\end{theorem}

\begin{proof}
Classical connection formula for Legendre functions of complex degree
\cite{Erdelyi}, specialized to the conical line. Verified at
$(\tau,x)=(0.6,3.0)$ and $(1.9,5.5)$ to relative error below
$1.3\times10^{-41}$.
\end{proof}

This is the reason the shadow transform behaves like an involution with
a fixed line rather than a generic reflection: on the principal series
it is a symmetry of the underlying Legendre equation, which depends on
$\nu$ only through $\nu(\nu+1)$.

% =====================================================================
\section{Exact prefactor and boundedness}
% =====================================================================

\begin{theorem}[Prefactor modulus]
\label{thm:prefactor}
For $h=\tfrac12(1+i\lambda)$, $\lambda>0$,
\begin{equation}
|c(\lambda)|^2
=\frac{4\,|\Gamma(1+i\lambda)|^2}{|\Gamma(\tfrac12+\tfrac{i\lambda}{2})|^4}
=\frac{2\lambda}{\pi}\coth\!\left(\frac{\pi\lambda}{2}\right).
\label{eq:prefactor}
\end{equation}
\end{theorem}

\begin{proof}
Reflection gives $|\Gamma(1+i\lambda)|^2=\pi\lambda/\sinh(\pi\lambda)$
and $|\Gamma(\tfrac12+\tfrac{i\lambda}{2})|^2=\pi/\cosh(\pi\lambda/2)$;
then use $\sinh(\pi\lambda)=2\sinh(\pi\lambda/2)\cosh(\pi\lambda/2)$.
Verified to $10^{-31}$ at three values of $\lambda$.
\end{proof}

\begin{theorem}[Boundedness on physical kinematics]
\label{thm:tempered}
On the physical slice $z=e^{i\varphi}$, $\bar z=z^*$, with $\varphi$ in
a compact subset of $(0,2\pi)$ avoiding $\varphi=0$, the scalar block
$g_{1+i\lambda}(z,\bar z)=k_h(z)k_h(\bar z)$ obeys
\begin{equation}
\bigl|g_{1+i\lambda}(z,\bar z)\bigr|
=|1-z|^{-1/2}\bigl(1+O(\lambda^{-1})\bigr),
\qquad \lambda\to+\infty .
\label{eq:tempered}
\end{equation}
No exponential growth survives: the rates $\pm\tfrac12\mathrm{Im}\,\zeta$
of the two chiral factors are conjugate-odd and cancel.
\end{theorem}

\begin{proof}
Uniform large-degree asymptotics of $Q_\nu$ at fixed argument, applied
to each chiral factor, combined with \eqref{eq:prefactor} in the form
$|c|^2=(2\lambda/\pi)(1+O(e^{-\pi\lambda}))$. Checked numerically at
$\varphi=1.0$ and $\varphi=2.2$: the ratio of $|g|$ to $|1-z|^{-1/2}$
equals $1.0059$, $0.99989$, $0.99914$ at $\lambda=4,12,30$, and
$0.99836$, $0.99892$, $0.99899$ respectively.
\end{proof}

\begin{remark}
The exponentially small correction $O(e^{-\pi\lambda})$ sits in the
prefactor, not in the block. The earlier abstract described an
``$e^{-\pi\lambda}$ envelope'' for the block itself, which the numbers
above contradict: at fixed $z$ inside the disc, $|k_h(z)|$ approaches a
constant near $0.72$ as $\lambda$ grows. Temperedness, in the sense
that matters for spectral decomposition, is the boundedness statement
\eqref{eq:tempered}.
\end{remark}

% =====================================================================
\section{The first moment, in closed form}
% =====================================================================

Write $P(\lambda)=\pi\lambda/\sinh(\pi\lambda)$, which the companion
paper derives as the phase-space factor
$\Gamma(1+i\lambda)\Gamma(1-i\lambda)$ of the two-particle cut on the
shadow locus \cite{ToupinLoops}. Its Fourier pair is elementary:
\begin{equation}
P(\lambda)=\int_\R e^{i\lambda x}\,\frac{\dd x}{4\cosh^2(x/2)},
\qquad
\frac{1}{2\pi}\int_\R P(\lambda)\cos(\lambda y)\,\dd\lambda
=\frac{1}{4\cosh^2(y/2)} .
\label{eq:fourier}
\end{equation}

\begin{theorem}[First moment against the digamma]
\label{thm:moment}
\begin{equation}
\frac{1}{2\pi}\int_\R P(\lambda)\,
\mathrm{Re}\,\psi\!\left(\tfrac12+\tfrac{i\lambda}{2}\right)\dd\lambda
\;=\;\frac18+\frac14\,\psi\!\left(\tfrac12\right).
\label{eq:moment}
\end{equation}
\end{theorem}

\begin{proof}
Subtract the constant. By \eqref{eq:fourier},
$\tfrac{1}{2\pi}\int_\R P=\tfrac14$, which produces the term
$\tfrac14\psi(\tfrac12)$. For the remainder use the Gauss integral
representation
$\psi(s)=-\gamma+\int_0^\infty\frac{e^{-t}-e^{-st}}{1-e^{-t}}\dd t$,
which gives
\begin{equation}
\mathrm{Re}\,\psi\!\left(\tfrac12+\tfrac{i\lambda}{2}\right)-\psi\!\left(\tfrac12\right)
=\int_0^\infty \frac{e^{-t/2}\bigl[1-\cos(\lambda t/2)\bigr]}{1-e^{-t}}\,\dd t .
\end{equation}
Exchange the order of integration (the integrand is positive) and apply
\eqref{eq:fourier} at $y=t/2$, together with
$e^{-t/2}/(1-e^{-t})=1/(2\sinh(t/2))$:
\begin{equation}
\frac{1}{2\pi}\int_\R P\,\bigl[\mathrm{Re}\,\psi-\psi(\tfrac12)\bigr]
=\int_0^\infty\frac{1}{2\sinh(t/2)}\cdot
\frac{1}{4}\Bigl[1-\operatorname{sech}^2(t/4)\Bigr]\dd t .
\end{equation}
Substituting $u=t/4$ and using
$1/(\sinh u\cosh u)=\operatorname{sech}^2u/\tanh u$, the integrand
collapses to $\tanh u\,\operatorname{sech}^2u$, so the remainder equals
\begin{equation}
\frac14\int_0^\infty \tanh u\,\operatorname{sech}^2u\,\dd u
=\frac14\left[\frac{\tanh^2u}{2}\right]_0^\infty=\frac18 .
\end{equation}
Each step was checked numerically to at least 25 digits, and the final
value matches the direct quadrature of \eqref{eq:moment} exactly at
working precision.
\end{proof}

\begin{remark}[The agreement of the two values of $1/8$]
\label{rem:coincidence}
The ladder moment of \cite{ToupinSpectral} is
$\mathcal M_1=\tfrac{1}{2\pi}\int_0^\infty P\,\dd\lambda=\tfrac18$,
taken over the half line, while the $\tfrac18$ in \eqref{eq:moment}
comes from an integral over the full line. The earlier draft read the
agreement as evidence that the one-loop measure appears inside the
digamma moment. Section \ref{sec:resolved} shows that the two numbers
really are the same quantity, and identifies what makes them so. The
reason has nothing to do with loops.
\end{remark}

% =====================================================================
\section{Why the two values agree}
\label{sec:resolved}
% =====================================================================

Write $\hat p$ for the Fourier partner of $P$,
\begin{equation}
\hat p(x)=\frac{1}{2\pi}\int_\R P(\lambda)e^{-i\lambda x}\dd\lambda
=\frac{1}{4\cosh^2(x/2)},
\qquad \hat p(0)=\frac14,\quad \hat p(\infty)=0 .
\label{eq:phat}
\end{equation}
Both quantities in question are integrals of $\hat p$ against the
kernel $1/\sinh$, so the question is what $\hat p$ does to that kernel.

\begin{theorem}[Weight-shift relations and the resulting differential equation]
\label{thm:ode}
Let $P(\lambda)=\Gamma(1+i\lambda)\Gamma(1-i\lambda)$. Then
\begin{equation}
P(\lambda\mp i)=\frac{1\pm i\lambda}{\mp i\lambda}\,P(\lambda),
\qquad
-i(\lambda-i)P(\lambda-i)+i(\lambda+i)P(\lambda+i)=4P(\lambda),
\label{eq:shift}
\end{equation}
and $\hat p$ satisfies, for all $x$, the first-order equation
\begin{equation}
\tfrac12\sinh(x)\,\hat p\,'(x)=\hat p(x)-\hat p(0).
\label{eq:ode}
\end{equation}
Equation \eqref{eq:ode} together with $\hat p(0)=\tfrac14$ and
$\hat p(\infty)=0$ determines $\hat p$ uniquely.
\end{theorem}

\begin{proof}
The first relation in \eqref{eq:shift} is the recursion
$\Gamma(z+1)=z\Gamma(z)$ applied to each factor; the second follows by
combining the two signs. Since $\lambda\mp i$ corresponds to
$\Delta\mapsto\Delta\pm1$ under $\Delta=1+i\lambda$, these are the
integer weight shifts on the principal series. For \eqref{eq:ode},
multiplication by $e^{\pm x}$ in the variable $x$ is the shift
$\lambda\mapsto\lambda\mp i$ on the transform side, and
differentiation is multiplication by $-i\lambda$; so \eqref{eq:shift}
transforms into $\tfrac12\sinh(x)\hat p\,'=\hat p-\hat p(0)$, the
constant arising because the shifted combination is integrable while
$\hat p(0)$ alone is not. Direct check: with
$\hat p=\tfrac14\operatorname{sech}^2(x/2)$ one has
$\hat p\,'=-\tfrac14\operatorname{sech}^2(x/2)\tanh(x/2)$, so the left
side is $\tfrac14\tanh^2(x/2)$, which is $\hat p(0)-\hat p(x)$ up to
sign. Verified at four values of $x$ to absolute error below
$1.3\times10^{-32}$. For uniqueness, separate variables:
$\dd\hat p/(\hat p-\tfrac14)=2\,\dd x/\sinh x$ integrates to
$\hat p-\tfrac14=C\tanh^2(x/2)$, and $\hat p(\infty)=0$ forces
$C=-\tfrac14$.
\end{proof}

\begin{theorem}[The two numbers are one boundary value]
\label{thm:resolved}
\begin{equation}
\frac{1}{2\pi}\int_\R P(\lambda)\Bigl[\mathrm{Re}\,\psi(\tfrac12+\tfrac{i\lambda}{2})-\psi(\tfrac12)\Bigr]\dd\lambda
=\int_0^\infty\frac{\hat p(0)-\hat p(x)}{\sinh x}\,\dd x
=-\frac12\int_0^\infty \hat p\,'(x)\,\dd x
=\frac{\hat p(0)}{2}
=\mathcal M_1 .
\label{eq:resolved}
\end{equation}
\end{theorem}

\begin{proof}
The first equality is the Fourier step of Theorem \ref{thm:moment},
using
$\mathrm{Re}\,\psi(\tfrac12+\tfrac{i\lambda}{2})-\psi(\tfrac12)
=\int_0^\infty(1-\cos\lambda x)/\sinh x\,\dd x$.
The second is \eqref{eq:ode} divided by $\sinh x$. The third is the
fundamental theorem of calculus with $\hat p(\infty)=0$. The last holds
because $\mathcal M_1=\tfrac{1}{2\pi}\int_0^\infty P\,\dd\lambda
=\tfrac12\cdot\tfrac{1}{2\pi}\int_\R P\,\dd\lambda=\tfrac12\hat p(0)$,
$P$ being even.
\end{proof}

So the agreement is a theorem, and the mechanism is visible: the
integrand of the digamma moment is an exact derivative, and what
survives is the value of $\hat p$ at the origin. That same value, halved,
is $\mathcal M_1$. Two integrals that look unrelated are two ways of
evaluating $\hat p(0)/2$, and the equation that makes the first one
collapse is the weight-shift relation \eqref{eq:shift} of the principal
series. Nothing about loop amplitudes enters, which is the sense in
which the earlier reading was wrong: the $1/8$ is a property of the
Gamma factors, and $\mathcal M_1$ inherits it.

\begin{proposition}[The transform is a thermal window]
\label{prop:fermi}
$\hat p(x)=-n_F'(x)$, where $n_F(x)=(e^x+1)^{-1}$ is the Fermi--Dirac
distribution at unit temperature in the variable $x$. Under the
Bisognano--Wichmann reading of \cite{ToupinSpectral}, $x$ is boost
rapidity and the temperature is $\beta=2\pi$ in boost time.
\end{proposition}

\begin{remark}[Which archimedean factor each object comes from]
$P=\Gamma(1+i\lambda)\Gamma(1-i\lambda)$ is built from Gamma factors of
the type attached to a complex place, while
$\mathrm{Re}\,\psi(\tfrac12+\tfrac{i\lambda}{2})$ is the logarithmic
derivative of $\Gamma(s/2)$ at $s=\Delta=1+i\lambda$, the type attached
to a real place. Theorem \ref{thm:resolved} therefore pairs the two.
We record the observation and stop there. No arithmetic consequence
follows from it, and readers of the earlier drafts should note the
vacuous zeta-zero check discussed in Sec.~\ref{sec:bridge} as a
cautionary case.
\end{remark}

\subsection{Higher moments}

The mechanism above extends into a recursion. Put
$f_k(y):=\tfrac{1}{2\pi}\int_\R P(\lambda)D(\lambda)^k e^{-i\lambda y}\dd\lambda$
with $D(\lambda)=\mathrm{Re}\,\psi(\tfrac12+\tfrac{i\lambda}{2})-\psi(\tfrac12)$,
so $f_0=\hat p$ and $A_k:=f_k(0)=\tfrac{1}{2\pi}\int_\R P D^k\dd\lambda$.
Define $T[f]:=\int_0^\infty[f(0)-f(y)]/\sinh y\,\dd y$. Then
\begin{equation}
A_{k+1}=T[f_k],
\qquad
T[f]=\tfrac12 f(0)-\int_0^\infty\frac{L[f](x)}{\sinh x}\dd x,
\qquad
L[f]:=\tfrac12\sinh(x)f'(x)-f(x)+f(0),
\label{eq:recursion}
\end{equation}
and Theorem \ref{thm:ode} says exactly that $L[\hat p]=0$, which is why
the first step closes. Equivalently $A_k=(-1)^k(\mathcal L^k\hat p)(0)$
for the operator
$\mathcal L f(x)=\int_0^\infty[f(x+y)+f(x-y)-2f(x)]\,\dd y/(2\sinh y)$,
the generator of the symmetric L\'evy process with jump density
$1/(2\sinh|y|)$, whose characteristic exponent is $D$ itself.
For $k\ge1$ the defect $L[f_k]$ does not vanish:
in the transform variable it involves $\mathrm{Re}\,\psi(1+i\lambda)$,
through the shifts
\begin{equation}
D(\lambda-i)-D(\lambda+i)=-\frac{2i}{\lambda},
\qquad
D(\lambda-i)+D(\lambda+i)=2\bigl[\mathrm{Re}\,\psi(\tfrac{i\lambda}{2})-\psi(\tfrac12)\bigr],
\label{eq:Dshift}
\end{equation}
both verified to 15 digits. The moments themselves are striking, and we
report them with their proof status marked.

\begin{center}
\begin{tabular}{llll}
\toprule
$k$ & $A_k$ & Status & Check \\
\midrule
0 & $1/4$ & proved & exact \\
1 & $1/8$ & proved (Thm.~\ref{thm:resolved}) & exact \\
2 & $1/8$ & \textbf{open} & 40 digits \\
3 & $\tfrac12\log 2-\tfrac{3}{16}$ & \textbf{open} & 30 digits \\
4 & $0.2315253305644009308\ldots$ & no closed form found & n/a \\
\bottomrule
\end{tabular}
\end{center}

\begin{theorem}[The second moment]
\label{thm:A2}
$A_2=\tfrac18=A_1$.
\end{theorem}

\begin{proof}
Substitute $\varphi=\tanh(x/2)$, under which
$\dd x/\sinh x=\dd\varphi/\varphi$ and
$a(x):=\hat p(0)-\hat p(x)=\tfrac14\varphi^2$. Expanding the two cosine
factors,
\begin{equation}
\frac{1}{2\pi}\int_\R P(\lambda)(1-\cos\lambda x)(1-\cos\lambda y)\dd\lambda
=a(x)+a(y)-\tfrac12 a(x+y)-\tfrac12 a(x-y),
\end{equation}
so with $\varphi,\psi$ the images of $x,y$ and $s=\varphi\psi$, the
addition rule $\tanh\frac{x\pm y}{2}=(\varphi\pm\psi)/(1\pm s)$ gives
\begin{equation}
A_2=\frac14\int_0^1\!\!\int_0^1
\Bigl[\varphi^2+\psi^2-\frac{(\varphi+\psi)^2}{2(1+s)^2}
-\frac{(\varphi-\psi)^2}{2(1-s)^2}\Bigr]\frac{\dd\varphi\,\dd\psi}{\varphi\psi} .
\end{equation}
The bracket collapses to
$s^2[(\varphi^2+\psi^2)(s^2-3)+4]/(1-s^2)^2$. Setting $u=\varphi^2$,
$v=\psi^2$ turns $\varphi\psi\,\dd\varphi\,\dd\psi$ into
$\tfrac14\dd u\,\dd v$, so
\begin{equation}
A_2=\frac{1}{16}\int_0^1\!\!\int_0^1
\frac{(u+v)(uv-3)+4}{(1-uv)^2}\,\dd u\,\dd v .
\label{eq:A2rational}
\end{equation}
Expand $(1-uv)^{-2}=\sum_{n\ge0}(n+1)(uv)^n$ and integrate term by
term, using $\int_0^1\!\int_0^1 u^av^b(uv)^n=[(a{+}n{+}1)(b{+}n{+}1)]^{-1}$.
With $m=n+1$ the summand telescopes:
\begin{equation}
\sum_{m\ge1}\Bigl[\frac{2m}{(m+1)(m+2)}-\frac{6}{m+1}+\frac4m\Bigr]
=\sum_{m\ge1}\Bigl[\frac4m-\frac{8}{m+1}+\frac{4}{m+2}\Bigr]=2 ,
\end{equation}
the coefficients summing to zero so that the harmonic parts cancel and
only the boundary terms $8-6$ survive. Hence $A_2=2/16=\tfrac18$.
Quadrature of \eqref{eq:A2rational} returns $2.0$ at working precision.
\end{proof}

A second route reaches the same value and produces a tool needed below.
Applying $T$ to $\hat p$ in closed form gives the generator image
\begin{equation}
f_1(x)=\frac{1-c}{8}\Bigl[2+\frac{(1+c)\log(1-c)}{c}\Bigr],
\qquad c=\tanh^2(x/2),
\label{eq:f1closed}
\end{equation}
verified against direct quadrature of $-\mathcal L\hat p$ at four points
to relative error below $3.1\times10^{-30}$, with $f_1(0)=\tfrac18=A_1$.
Then $A_2=T[f_1]=\tfrac{1}{16}\int_0^1[-1+2c-(1-c^2)c^{-1}\log(1-c)]\dd c/c$,
whose integrand has the cancellation-free form
$\tfrac52+c^{-2}[-\log(1-c)-c-\tfrac{c^2}{2}]+\log(1-c)$ and integrates
to $\tfrac52+\tfrac12-1=2$.

\subsection{Status of the remaining moments}

Normalizing $4P(\lambda)\dd\lambda/2\pi$ to a probability measure, $D$
has mean $A_1/A_0=\tfrac12$ and variance
$A_2/A_0-(A_1/A_0)^2=\tfrac14$ by Theorem \ref{thm:A2}, so its standard
deviation equals its mean.

For $A_3$ the same machinery applies and the computation is longer. Two
ingredients are in place. First, \eqref{eq:f1closed} reduces $A_3$ to a
two-dimensional integral of elementary functions,
\begin{equation}
A_3=\int_0^\infty\!\!\int_0^\infty
\frac{b(x)+b(y)-\tfrac12 b(x+y)-\tfrac12 b(x-y)}{\sinh x\,\sinh y}\,\dd x\,\dd y,
\qquad b:=f_1(0)-f_1 ,
\label{eq:A3reduction}
\end{equation}
which quadrature confirms against
$\tfrac12\log2-\tfrac{3}{16}$ to eleven digits. Second, the logarithm in
\eqref{eq:f1closed} splits under addition, since
\begin{equation}
1-\tanh^2\!\Bigl(\frac{x+y}{2}\Bigr)
=\frac{(1-\varphi^2)(1-\psi^2)}{(1+\varphi\psi)^2},
\end{equation}
so $\log(1-c)$ is additive up to $\log(1+\varphi\psi)$; the alternating
denominator is the natural source of the $\log 2$. We have not carried
the evaluation through, and record $A_3$ as open.

\begin{center}
\begin{tabular}{llll}
\toprule
$k$ & $A_k$ & Status & Check \\
\midrule
0 & $1/4$ & proved & exact \\
1 & $1/8$ & proved (Thm.~\ref{thm:resolved}) & exact \\
2 & $1/8$ & proved (Thm.~\ref{thm:A2}) & exact \\
3 & $\tfrac12\log 2-\tfrac{3}{16}$ & open, reduced to \eqref{eq:A3reduction} & 30 digits \\
4 & $0.2315253305644009308\ldots$ & no closed form found & n/a \\
\bottomrule
\end{tabular}
\end{center}

\begin{thebibliography}{9}
\bibitem{ToupinLoops} D.~Toupin, ``Loop amplitudes from unitarity cuts
in celestial Mellin space,'' v19, 2026.
\bibitem{ToupinSpectral} D.~Toupin, ``The spectral weight
$\pi\lambda/\sinh\pi\lambda$: derivation, modular form, and exact
moments,'' v2, 2026.
\bibitem{Erdelyi} A.~Erd\'elyi et al., \textit{Higher Transcendental
Functions}, Vol.~1, McGraw--Hill, 1953.
\bibitem{LamShao2018} H.~T.~Lam, S.-H.~Shao, Phys.\ Rev.\ D98 (2018)
025020 [arXiv:1711.06138].
\bibitem{Optical2024} ``The celestial optical theorem,''
arXiv:2404.18898.
\end{thebibliography}

\end{document}
